% Candidate introduction/related-work text, not included by the main document.
% Uses references.bib and Literature_Review_Additions.bib.
% Apply Literature_Review_Replacements.bib by replacing matching entries,
% not by loading a second database containing the same keys.
% The final paragraph intentionally describes scope without claiming priority
% or asserting completion of the proof obligations in Literature_Review.md.

Robust control of distributed-parameter systems requires methods that account
for both spatial dynamics and uncertainty. Established approaches include PDE
backstepping, direct Lyapunov methods for uncertain parabolic systems, and
semidefinite programming for PDE stability
analysis~\cite{krstic2008boundary,fridman2012sampled,valmorbida2016pdesdp}.
The partial integral equation (PIE) representation provides an algebraic
description of a class of PDEs through bounded partial integral (PI)
operators~\cite{peet2021pierepresentation,shivakumar2024GDPE}.
Its dual representation enables nominal state-feedback synthesis using
operator inequalities~\cite{shivaknew}. Extending this construction to
uncertain systems requires uncertainty descriptions compatible with both the
PIE representation and controller synthesis.

Integral quadratic constraints (IQCs) provide a framework for describing
uncertainty and nonlinear input--output behavior~\cite{megretski1997system}.
Their dissipation interpretation connects these descriptions to storage
inequalities~\cite{willems1972dissipative1,6915700}.
PI multipliers have been used for robust ODE--PDE
analysis~\cite{9303892}, and hard IQCs with infinite-dimensional channels
admit sufficient PIE-based KYP tests~\cite{10156335}.
More recently, an IQC-based PIE framework has addressed structured robustness,
performance, and observer-gain
synthesis~\cite{lenssen2025muanalysissynthesisframeworkpartial}.
These results motivate a corresponding treatment of robust feedback
controller synthesis.

IQC duality also has established finite-dimensional synthesis applications.
Dual constraints enable robust disturbance-feedforward
design~\cite{kose2009robust}, while dual filters and storage transformations
support convex estimator and feedforward synthesis for gridded LPV
systems~\cite{venkataraman}.
Convexity depends on the plant and controller
structure~\cite{scherer2009robust}.
For more general dynamic multipliers, synthesis methods commonly alternate
between controller design and multiplier
analysis~\cite{veenman2014iqcsynthesis,wang2016robustiqc}.
Dynamic-IQC feedback constructions also exist for uncertain systems with
state delay~\cite{7353149}. These distinctions are relevant when comparing a
convex synthesis condition with a joint controller--multiplier optimization.

The passage between frequency-domain multipliers and finite-horizon
inequalities requires particular care. Factorizations preserving graph and
inverse-graph conditions are studied through doubly-hard
factorizations~\cite{carrasco2018equivalence}; finite-horizon IQCs with terminal
costs provide a complementary connection to
dissipativity~\cite{scherer2018dynamicdissipation}.
Additional storage associated with the multiplier can also reduce
conservatism~\cite{magicalhiddenenergy,hu}.
These results distinguish admissible finite-horizon certificates from
frequency-domain factorization identities alone.

This paper studies robust state-feedback synthesis for uncertain systems
admitting a PIE representation through duality of dynamically filtered IQC
conditions. The construction addresses the placement of the reconstruction
operator in the synthesis inequality and the recovery of controllers from
dual operator variables. The synthesis problem is considered for fixed
admissible filter data, with explicit causal inverse and realization
assumptions. Static factors lead to state feedback, while dynamic factors
introduce controller states driven by specified available signals. The
uncertainty classes and measurement assumptions are stated separately for
these two realizations.
